\section{Related Work}\label{sec:related}

%\subsection{Data Shifts}
%Data shifts pose a fundamental challenge in machine learning, broadly categorized into \emph{covariate shift}, \emph{label shift}, and \emph{concept shift}, etc.~\cite{quionero2009dataset}. To systematically examine data shifts, we categorize them into two main types \cite{10.5555/3692070.3692237}:
%(i) Data Corruptions: alterations in the distribution’s support, i.e., changes in the observed data instances themselves, often arising from measurement errors, environmental disturbances, or adversarial manipulations. Examples include background noise affecting speech recognition systems~\cite{kinoshita2020deep}, numerical rounding errors in financial computations~\cite{li2015rounding}, adversarial modifications in image classification tasks~\cite{goodfellow2020explaining}, and distortions in decision-making processes caused by Goodhart’s law in public aid allocation~\cite{camacho2011goodhart};
%(ii) Subpopulation Shifts: changes to the distribution’s density or mass function while its support remains unchanged. Notable examples include demographic changes influencing recommendation systems~\cite{blodgett2016demographic,sapiezynski2017algorithmic}, temporal shifts affecting diagnostic accuracy in healthcare~\cite{pasterkamp2017moving}, and spatial distribution changes disrupting wildlife conservation models~\cite{beery2021iwildcam}. By proactively modeling these shifts within the training dataset, rather than reacting to them post-deployment, our method builds robustness into CF generation directly.

%\subsection{Robust CFs}

%Ensuring robustness under various perturbations remains a key research challenge. Prior work has explored robustness from multiple perspectives, including robustness against model changes, model multiplicity, noisy execution, and input variations. Among these, our work aligns most closely with robustness against model changes and model multiplicity. Therefore, we provide a more detailed discussion on them. For a comprehensive overview of other robustness methods, we refer readers to the relevant survey paper \cite{ijcai2024p894}.

%\subsubsection{Robustness against Model Changes}
Model parameter updates typically arise from retraining on new data~\cite{rawal2020algorithmic}. Existing approaches constrain such changes either \emph{implicitly}, without modeling their source, or \emph{explicitly}, by model plausible parameter shifts .

\textbf{Implicit constraints.}
The methods optimize CFs for high classification certainty~\cite{dutta2022robust,Hamman2023Robust,DBLP:conf/aaai/JiangL0T23}, augment training data with generated CFs~\cite{ferrario2022robustness}, or perturb the covariance of data sampled near the decision boundary to train a linear surrogate~\cite{DBLP:journals/ior/BuiNYN25}. \citet{xu2024generally} relax the naturally-occurring model change condition of \citet{Hamman2023Robust} to a subgaussian bound. Two works certify robustness a posteriori under a Bayesian view over a user-specified pool of retrained models: \citet{stkepka2025counterfactual} move a base CF into a region meeting a specified robustness level, while \citet{duell2026provably} optimize on the generation side for high confidence and low variance in the CF's posterior predictive distribution, certifying graceful degradation via a KL-divergence bound.


\textbf{Explicit constraints.}
\citet{upadhyay2021towards} introduce ROAR, %the notion of plausible model shifts and 
a minimax objective generating CFs that minimize the worst-case plausible shift under a linear surrogate %of the (possibly nonlinear) 
model; \citet{kayastha2024learning} give a provably optimal solver for any generalized linear model, and \citet{kyaw2025optimal} extend this to general $L_p$-bounded neighborhoods. \citet{DBLP:conf/iclr/NguyenBN23} similarly hedge against plausible parameter shifts via a min-max problem, where the max term is the Gelbrich distance to an ambiguity set of future model-parameter distributions, solved by gradient descent. \citet{bui2022counterfactual} instead bound the joint validity probability of a set of CFs. AP$\Delta$S~\cite{Marzari2024Rigorous,marzari2025probabilistically} samples plausible model shifts directly and shows exact robustness decisions are NP-hard. \citet{DBLP:conf/acml/JiangLL0T23,jiang2024interval} propagate interval-valued weights and biases through the network to bound hidden-unit activations under bounded parameter perturbations.
%A limitation shared by these explicit, parameter-space approaches is that they assume model changes stem from minor data shifts without modeling the data shift itself. 
These parameter-space approaches assume model changes stem from data shifts without modeling the shift itself.
\citet{nguyen2022robust} address this by adversarially perturbing the data via a Gaussian-mixture ambiguity set over Wasserstein distance on the mean-covariance space.
Two works come closest to a population view: \citet{bui2022counterfactual} jointly optimize a small per-instance CF set for joint validity, while \citet{nguyen2022robust}, despite modeling the data shift, still certifies recourse one input at a time. Neither treats the population of generated CFs as the robustness object, the gap our method closes.

A related line of work study robustness to \emph{model multiplicity}~\cite{breiman2001statistical,Marx2020Predictive,Black2022Selective,Rudin2019Stop,coston2021characterizing,pawelczyk2020counterfactual,turbal2026ellice}, where distinct near-optimal models yield inconsistent CFs, with no assumption on how or why models differ. See~\cite{ijcai2024p894} for a broader robustness survey.

%Most existing research on CF robustness focuses on ensuring validity under model parameter updates, which typically occur within the model's original structural framework. Such modifications arise when models are retrained on data from slightly altered distributions \cite{rawal2020algorithmic}, aligning with our study's framework. Recent works have explored various ways to define the boundaries of model changes. For instance, \cite{upadhyay2021towards} introduced the concept of plausible model variations, constraining parameter updates to perturbations within a predefined magnitude. Formally, given an input $x$, model $f$, and CF $x'$, their framework ensures that CFs remain valid under bounded parameter shifts, requiring both the original model ff and the perturbed model $f'$ to assign the same classification to $x'$. This fundamental idea has been extended in \cite{Black2022Consistent,Jiang2023Provably,Jiang2023Formalising,Marzari2024Rigorous}, all maintaining the core requirement that CFs remain consistent across model updates. In contrast, \cite{Hamman2023Robust} expanded the scope beyond artificially bounded perturbations, considering naturally occurring model shifts in neural networks. Their approach eliminates explicit magnitude constraints, allowing theoretically unbounded parameter changes while ensuring local consistency in model behavior near the explained input. While these methodologies vary in how they define acceptable model changes—some imposing strict constraints while others allowing greater flexibility—there is a shared recognition that model instability undermines the reliability of explanations. This issue is particularly evident in high-stakes scenarios, such as fluctuating loan approval outcomes across model updates. A key limitation of existing approaches is their assumption that model changes stem from small data shifts, yet they do not explicitly analyze these shifts in depth. Since future data distributions remain uncertain, accounting for such shifts makes the problem more complex. In contrast, our method takes a proactive approach by modeling and anticipating the most probable data shifts, ensuring greater robustness in CFs.





%Model multiplicity refers to the phenomenon where multiple distinct models achieve comparable predictive performance for the same task \cite{breiman2001statistical,Marx2020Predictive}. While one common approach aggregates predictions from these models into a unified output \cite{Black2022Selective}, recent studies indicate that such ensemble models may still exhibit inconsistent behavior — both in terms of predictions and their explanations — when applied to individual instances \cite{Rudin2019Stop,coston2021characterizing}. This inconsistency has significant implications for CFs: a CF generated for an aggregated prediction may not remain valid across all constituent models, potentially leading to contradictions and undermining user trust in explanation consistency. To address this, \cite{pawelczyk2020counterfactual} anchor explanations to a specific model's predictions and ensure CF validity across potential classifier updates, mirroring the goals of robustness against model changes. Unlike that line of work, model multiplicity does not assume updates are driven by shifts in data distribution, and places no constraints on the architecture or training methodology of either the base model generating CFs or the alternative models considered.
